\chapter{Hashimoto's Disease}

\section*{Definition and Overview}
Hashimoto's disease, also known as Hashimoto thyroiditis or chronic lymphocytic thyroiditis, is an autoimmune disorder in which the body's immune system mistakenly attacks the thyroid gland, causing inflammation and gradual destruction of thyroid tissue. It is the most common cause of an underactive thyroid (hypothyroidism). The condition develops slowly, often over many years, and many people have no symptoms initially. As thyroid function declines, symptoms such as fatigue, weight gain, cold intolerance, and constipation appear. Hashimoto's disease is easily diagnosed with blood tests and treated with daily thyroid hormone replacement, allowing people to live normally.

\section*{Epidemiology}
Hashimoto's disease is the most common autoimmune disease and the leading cause of hypothyroidism in Australia. It affects around 1 in 20 to 1 in 30 people, and it is five to ten times more common in women than in men. It most often begins between ages 30 and 50, and prevalence increases with age. It runs in families and is associated with other autoimmune conditions, including coeliac disease, type 1 diabetes, rheumatoid arthritis, and pernicious anaemia. Many people with the disease have a goitre at some stage.

\section*{Aetiology and Pathophysiology}
Hashimoto's disease results from an autoimmune attack on the thyroid gland.

\begin{itemize}
    \item \textbf{Autoantibodies:} the immune system produces antibodies against thyroid peroxidase (TPO) and thyroglobulin
    \item \textbf{Chronic inflammation:} lymphocytes infiltrate the gland, gradually destroying the thyroid follicles that produce hormone
    \item \textbf{Hypothyroidism:} as the gland is destroyed, thyroid hormone production falls and TSH (thyroid-stimulating hormone) rises
    \item \textbf{Goitre:} the gland often enlarges early in the disease in response to inflammation and stimulation
    \item \textbf{Genetic susceptibility:} certain genes predispose to the condition
    \item \textbf{Environmental triggers:} infections, stress, pregnancy, and excess iodine can trigger onset
    \item \textbf{Progression:} most people eventually develop permanent hypothyroidism
\end{itemize}

\section*{Clinical Features}
Symptoms develop gradually and are often dismissed as normal ageing or stress.

\begin{itemize}
    \item Fatigue, tiredness, and weakness
    \item Weight gain and difficulty losing weight
    \item Feeling cold and cold intolerance
    \item Constipation
    \item Dry skin, dry hair, and hair loss
    \item Depression, low mood, and brain fog
    \item Muscle aches, joint pain, and stiffness
    \item Heavy or irregular periods
    \item A goitre: swelling at the front of the neck, sometimes with a feeling of pressure
    \item Hoarseness and puffy face in advanced cases
    \item Slow heart rate
\end{itemize}

\section*{Differential Diagnosis}
Other causes of hypothyroidism and similar symptoms should be considered.

\begin{itemize}
    \item Iodine deficiency, which is uncommon in Australia since iodine fortification
    \item Postpartum thyroiditis
    \item Subacute thyroiditis, a painful viral inflammation
    \item Thyroid surgery or radioactive iodine treatment causing hypothyroidism
    \item Medicines affecting thyroid function, including lithium and amiodarone
    \item Other autoimmune diseases causing fatigue
    \item Anaemia, depression, and chronic fatigue
    \item Goitre from other causes, including nodules and cancer
\end{itemize}

\section*{When to Seek Medical Care}
See a doctor for persistent fatigue, weight gain, cold intolerance, or other symptoms of hypothyroidism, and for a neck swelling. Thyroid testing is routine and simple.

\textbf{Call 000 (or local emergency services) for:}
\begin{itemize}
    \item Severe confusion, low body temperature, or slowed breathing in a person with hypothyroidism, which may indicate myxoedema coma
    \item Chest pain, severe breathlessness, or an irregular heartbeat
    \item A rapidly enlarging neck lump with difficulty breathing or swallowing
\end{itemize}

\section*{Diagnosis and Investigations}
Hashimoto's disease is confirmed with blood tests.

\begin{itemize}
    \item \textbf{Thyroid function tests:} a high TSH with a low free T4 confirms hypothyroidism
    \item \textbf{Thyroid antibodies:} high TPO antibodies confirm the autoimmune cause
    \item \textbf{Ultrasound:} shows a characteristic gland appearance and checks for nodules
    \item \textbf{Liver function and cholesterol:} may be checked, as hypothyroidism affects these
    \item \textbf{Other autoimmune screening:} depending on symptoms and family history
    \item \textbf{Monitoring:} regular thyroid function tests during treatment
\end{itemize}

\section*{Management}
Treatment is safe, simple, and lifelong.

\begin{itemize}
    \item \textbf{Levothyroxine:} daily thyroid hormone replacement restores normal hormone levels
    \item \textbf{Dose individualisation:} the dose is tailored to blood test results
    \item \textbf{Regular monitoring:} thyroid function is checked until stable, then annually
    \item \textbf{Correct timing:} levothyroxine is taken on an empty stomach, usually 30--60 minutes before breakfast
    \item \textbf{Subclinical disease:} mild elevations of TSH may be treated depending on symptoms and risk factors
    \item \textbf{Pregnancy:} doses often increase, and close monitoring is essential
    \item \textbf{Goitre monitoring:} the goitre is observed, and hormone treatment may shrink it
    \item \textbf{Healthy lifestyle:} balanced diet, exercise, and adequate iodine
\end{itemize}

\section*{Complications}
\begin{itemize}
    \item Untreated hypothyroidism: heart disease, high cholesterol, and reduced quality of life
    \item Myxoedema coma, a rare but life-threatening complication
    \item Effects in pregnancy, including increased risk of miscarriage and effects on the baby's development if untreated
    \item Goitre with pressure symptoms
    \item Increased risk of other autoimmune conditions
    \item Rarely, thyroid lymphoma
    \item Consequences of over- or undertreatment with levothyroxine
\end{itemize}

\section*{Prognosis and Outcomes}
The prognosis for Hashimoto's disease is excellent. With daily levothyroxine and regular monitoring, symptoms resolve, blood tests return to normal, and people enjoy normal health, energy, and life expectancy. Treatment is lifelong, but it is one of the simplest and safest of all long-term therapies. Early diagnosis prevents the complications of untreated hypothyroidism, and pregnancy outcomes are normal with good management.

\section*{Prevention and Self-Care}
\begin{itemize}
    \item Take levothyroxine daily, at the same time, on an empty stomach
    \item Do not stop or change the dose without medical advice
    \item Attend regular blood tests
    \item Avoid taking iron, calcium, or antacids at the same time as levothyroxine
    \item Eat a balanced diet; do not take iodine supplements unless advised
    \item Report new symptoms or pregnancy promptly
    \item Have family members with symptoms tested, as the condition runs in families
    \item Maintain a healthy lifestyle with exercise and stress management
\end{itemize}

\section*{Sources and Further Reading}
The following reputable sources provide further information for healthcare students and patients seeking reliable, current advice about Hashimoto's disease.

\begin{itemize}
    \item Australian Thyroid Foundation. \textit{Hashimoto's disease information.} https://www.thyroidfoundation.org.au/
    \item Healthdirect Australia. \textit{Hashimoto's disease.} https://www.healthdirect.gov.au/
    \item Jean Hailes for Women's Health. \textit{Thyroid conditions.} https://www.jeanhailes.org.au/
    \item American Thyroid Association. \textit{Hashimoto's thyroiditis.} https://www.thyroid.org/
    \item NHS. \textit{Underactive thyroid.} https://www.nhs.uk/
    \item National Institute of Diabetes and Digestive and Kidney Diseases. \textit{Hashimoto's disease.} https://www.niddk.nih.gov/
\end{itemize}
